\begin{itemize}
    \color{red}
    \item introduce the concept of 4D tracking
    \item motivate the need for 4D tracking
    \begin{itemize}
        \item HL-LHC pileup challenge
        \item potential for physics performance
        \item potential for CPU performance
        \item opportunity in ATLAS with Phase-II upgrade and beyond
    \end{itemize}
\end{itemize}

Tracking sensors with timing information are a recent development in high energy physics. The ATLAS detector will have a partial coverage of timing information with the Phase-II ugparde.

To study the impact of the additional time measurements on the tracking and physics performance, assumtions have to be made about the time resolution of the sensors and the amount of material in the detector. In the best case this involves a detailed simulation of the detector and a reconstruction with algorithms which fully utilize this new dimension.

The additional information needs to be processed not only in hardward but also software. It needs to be stored in the EDM and processed by the reconstruction algorithms.

\section{Tracking with time information in ACTS}

Tracking with time information is a long standing feature in ACTS. Time is part of the track parameterization and extrapolated in the propagator. If measurements carry time information, it can be accessed in the same way as other measurement dimensions.

While some algorithms do not need special treatment to utilize time information, others need to be adapted and configured correctly to make use of it. At the moment the track EDM does not make a difference if parameters are fitted or not. For this reason tracks with and without time information can only be distinguished by looking into the measurement content or taking a decision based on the variance of the time parameter. The variance will be either on the order of the resolution of the time measurements or practically the initial variance of the time parameter. This is a limitation of the current implementation which has to be addressed in the future.

\section{Spacepoints with time}

After clusters and measurements with time information are created, the information can be further propagated into spacepoints. Spacepoints are a basic building block for pattern recognition algortihms and can be used to cut combinatorics in the seeding stage.

During the work on this thesis spacepoints were extended to carry optional time information. This is a necessary step to further propagate this information to the seed track parameters which is then automatically picked up by the track finding.

\section{Seeding with time}

Seeding is the first step in the track finding process. It is responsible for creating track candidates from spacepoints. This is a combinatorial problem and cuts are used to reduce the number of candidates. Time information can be used to cut the combinatorics even further, therefore reducing the computational cost, but this requires having at least two spacepoints with time information.

The current implementation of the triplet seeding in ACTS will not make use of time information. There has been an early study to estimate the potential of time information in seeding. The results are promising and show a significant reduction in the number of fakes and in reduction of the computational cost.

During the work on this thesis time information has been included in the parameter estimation of triplet seeds. This is a necessary step to further propagate this information to the track finding. It is a very simple implementation for now which just picks up the time information from the bottom spacepoint if it is available. This could be further improved by taking the time information from any spacepoint, extrapolate it to the bottom spacepoint and average it to use it for the parameters.

\section{Track finding and fitting with time}

\section{Vertex finding with time}

\section{Vertex fitting with time}
